%%%%%%%%%%%%%%%%%%%%%%%%%%%%%%%%%%%%%%%%%
% Professional Newsletter Template
% LaTeX Template
% Version 1.0 (09/03/14)
%
% Created by:
% Bob Kerstetter (https://www.tug.org/texshowcase/) and extensively modified by:
% Vel (vel@latextemplates.com)
% 
% This template has been downloaded from:
% http://www.LaTeXTemplates.com
%
% License:
% CC BY-NC-SA 3.0 (http://creativecommons.org/licenses/by-nc-sa/3.0/)
%
%%%%%%%%%%%%%%%%%%%%%%%%%%%%%%%%%%%%%%%%%

\documentclass[9pt]{extarticle} % The default font size is 10pt; 11pt and 12pt are alternatives

%%%%%%%%%%%%%%%%%%%%%%%%%%%%%%%%%%%%%%%%%
% Professional Newsletter Template
% Structural Definitions File
% Version 1.0 (09/03/14)
%
% Created by:
% Vel (vel@latextemplates.com)
% 
% This file has been downloaded from:
% http://www.LaTeXTemplates.com
%
% License:
% CC BY-NC-SA 3.0 (http://creativecommons.org/licenses/by-nc-sa/3.0/)
%
%%%%%%%%%%%%%%%%%%%%%%%%%%%%%%%%%%%%%%%%%

%----------------------------------------------------------------------------------------
%	REQUIRED PACKAGES
%----------------------------------------------------------------------------------------

\usepackage{listings}
\usepackage{graphicx} % Required for including images
\usepackage{microtype} % Improved typography
\usepackage{multicol} % Used for the two-column layout of the document
\usepackage{booktabs} % Required for nice horizontal rules in tables
\usepackage{wrapfig} % Required for in-line images
\usepackage{float} % Required for forcing figures not to float with the [H] parameter
\usepackage[utf8]{inputenc}
\usepackage{fancyhdr}

%------------------------------------------------
% Fonts

\usepackage{charter} % Use the Charter font as the main document font
\usepackage{courier} % Use the Courier font for \texttt (monospaced) only
\usepackage[T1]{fontenc} % Use T1 font encoding

%------------------------------------------------
% List Separation

\usepackage{enumitem} % Required to customize the list environments
\setlist{noitemsep,nolistsep} % Remove spacing before, after and within lists for a compact look

%------------------------------------------------
% Figure and Table Caption Styles

\usepackage{caption} % Required for changing caption styles
\captionsetup[table]{labelfont={bf,sf},labelsep=period,justification=justified} % Specify the table caption style
\captionsetup[figure]{labelfont={sf,bf},labelsep=period,justification=justified, font=small} % Specify the figure caption style
\setlength{\abovecaptionskip}{10pt} % Whitespace above captions

%------------------------------------------------
% Spacing Between Paragraphs

\makeatletter
\usepackage{parskip}
\setlength{\parskip}{6pt}
\newcommand{\@minipagerestore}{\setlength{\parskip}{6pt}}
\makeatother

%----------------------------------------------------------------------------------------
%	PAGE MARGINS AND SPACINGS
%----------------------------------------------------------------------------------------

\textwidth = 7 in % Text width
\textheight = 10 in % Text height
\oddsidemargin = -18pt % Left side margin on odd pages
\evensidemargin = -18pt % Left side margin on even pages
\topmargin = -36pt % Top margin
\headheight = 0pt % Remove the header by setting its space to 0
\headsep = 0pt % Remove the space between the header and top of the page
\parskip = 2pt % Space between paragraph
\parindent = 0.0in % Paragraph indentation
\pagestyle{empty} % Disable page numbering

%----------------------------------------------------------------------------------------
%	COLORS
%----------------------------------------------------------------------------------------

\usepackage[dvipsnames,svgnames]{xcolor} % Required to specify custom colors

\definecolor{altncolor}{rgb}{.8,0,0} % Dark red
%\definecolor{altncolor}{rgb}{.2,.4,.8} % Dark blue
%\definecolor{altncolor}{rgb}{.84,.16,.16} % Red

\usepackage[colorlinks=true, linkcolor=altncolor, anchorcolor=altncolor, citecolor=altncolor, filecolor=altncolor, menucolor=altncolor, urlcolor=altncolor]{hyperref} % Use the color defined above for all links

%----------------------------------------------------------------------------------------
%	BOX STYLES
%----------------------------------------------------------------------------------------

\usepackage[framemethod=TikZ]{mdframed}% Required for creating boxes
\mdfdefinestyle{sidebar}{
    linecolor=black, % Outer line color
    outerlinewidth=0.5pt, % Outer line width
    roundcorner=0pt, % Amount of corner rounding
    innertopmargin=10pt, % Top margin
    innerbottommargin=10pt, % Bottom margin
    innerrightmargin=10pt, % Right margin
    innerleftmargin=10pt, % Left margin
    backgroundcolor=white, % Box background color
    frametitlebackgroundcolor=white, % Title background color
    frametitlerule=false, % Title rule - true or false
    frametitlerulecolor=white, % Title rule color
    frametitlerulewidth=0.5pt, % Title rule width
    frametitlefont=\Large, % Title heading font specification
    font=\small
}

\mdfdefinestyle{intextbox}{
    linecolor=black, % Outer line color
    outerlinewidth=0.5pt, % Outer line width
    roundcorner=10pt, % Amount of corner rounding
    innertopmargin=7pt, % Top margin
    innerbottommargin=7pt, % Bottom margin
    innerrightmargin=7pt, % Right margin
    innerleftmargin=7pt, % Left margin
    backgroundcolor=white, % Box background color
    frametitlebackgroundcolor=white, % Title background color
    frametitlerule=false, % Title rule - true or false
    frametitlerulecolor=white, % Title rule color
    frametitlerulewidth=0.5pt, % Title rule width
    frametitlefont=\Large % Title heading font specification
}

%----------------------------------------------------------------------------------------
%	HEADING STYLE
%----------------------------------------------------------------------------------------

\newcommand{\heading}[2]{ % Define the \heading command
\vspace{#2} % White space above the heading
{\begin{center}\Large\textbf{#1}\end{center}} % The heading style
\vspace{#2} % White space below the heading
}

\newcommand{\BackToContents}{\hyperlink{contents}{{\small Back to Contents}}} % Define a command for linking back to the contents of the newsletter % Include the document which specifies all packages and structural customizations for this template

\def\changemargin#1#2{\list{}{\rightmargin#2\leftmargin#1}\item[]}
\let\endchangemargin=\endlist

\begin{document}

%--------------------------------------------------------------------------------
% HEADER DETAILS
%--------------------------------------------------------------------------------

\pagestyle{fancy}
\fancyhf{}
\chead{segfault@csh.rit.edu}
\rhead{\today}
\lhead{Volume XLVIII Issue \#4}
\addtolength\footskip{-15px}
\cfoot{"I'm just imagining Derek in a diaper now" - Will Ziener-Dignazio (slackwill)}

%----------------------------------------------------------------------------------------
%	HEADER IMAGE
%----------------------------------------------------------------------------------------

\begin{figure}[H]
\centering\vspace{0.5cm}\includegraphics[width=0.8\linewidth]{imgs/segfault.png}
\end{figure}

%--------------------------------------------------------------------------------
% HEADER QUOTE
%--------------------------------------------------------------------------------

\vspace{-15px}
\begin{quote}
\centering
\textbf{\textit{Just do it already}}
\end{quote}
\vspace{10px}

%----------------------------------------------------------------------------------------
%	SIDEBAR - FIRST PAGE
%----------------------------------------------------------------------------------------

\vspace{-0.5cm}\begin{minipage}[t]{.35\linewidth} % Mini page taking up 35% of the actual page
\begin{mdframed}[style=sidebar,frametitle={}] % Sidebar box

%-----------------------------------------------------------

\hypertarget{contents}{\textbf{{\large This week on floor\ldots}}} % \hypertarget provides a label to reference using \hyperlink{label}{link text}
\begin{itemize}
\item \hyperlink{firstnews}{Sequential Consistency}
\item \hyperlink{secondnews}{Client-Centric Consistency Models}
\end{itemize}

\centerline {\rule{.75\linewidth}{.25pt}} % Horizontal line

%-----------------------------------------------------------

\textbf{Notable Upcoming Events:}
\begin{enumerate}[leftmargin=0.2cm]
\item \textbf{<EVENT NAME>} <DATE + TIME> \\
	<DESCRIPTION>
\\
\item \textbf{<EVENT NAME>} <DATE + TIME> \\
	<DESCRIPTION>
\\
\end{enumerate}

%-----------------------------------------------------------


\textbf{<TITLE>} \\
<SHORT SECTION>
\\

%-----------------------------------------------------------

\captionof*{table}{Voting Results}
\begin{tabular}{lcr}

Vote & Cost & Result \\
\midrule
<NAME> & \$<MONEY> & <STATUS> \\
\bottomrule
\end{tabular}

%-----------------------------------------------------------

\end{mdframed}
\end{minipage}\hfill % End the sidebar mini page 
%
%----------------------------------------------------------------------------------------
%	MAIN BODY - FIRST PAGE
%----------------------------------------------------------------------------------------
%
\begin{minipage}[t]{.61\linewidth} % Mini page taking up 61% of the actual page
\vspace{-0.4cm}
\hypertarget{firstnews}{\heading{Sequential Consistency}{6pt}}
Sequential consistency is an important data-centic consistency model, which was first defined by Lamport in the context of shared memory for multiprocessor systems. In general, a data store is said to be sequentially consistent when it satisfies the following condition:
\begin{changemargin}{0.6cm}{0.6cm} 
\textit{The result of any execution is the same as if the (read and write) operations by all process on the data store were executed in some sequential order and the operations of each individual process appear in this sequence in the order specified by its program.}
\end{changemargin}
What this definition means is that when processes run concurrently on different machines, any valid interleaving of read and write operations is acceptable behavior, but \textit{all processes see the same interleaving of operations.} In this context, a process "sees" writes from all processes but only "sees" its own reads. An example of sequential consistency is seen in Lamport's Part-time parliament paper and Diego's Raft paper.

%-----------------------------------------------------------

\hypertarget{secondnews}{\heading{Client-Centric Consistency Models}{6pt}}

An important assumption is that concurrent processes may be simultaneously updating a data store, and that it is necessary to provide consistency in the face of such concurrency. Many consistency models try and provide a system-wide consistent view, but this can be expensive and time consuming. An alternative approach is called client-centric consistency models, in which consistency requirements are moved to the client's view. We will go over some of these now. \\
\\
\textbf{Monotonic-Reads:} If a process reads the value of data item x, any successive read operation on x by that process will always return the same value of a more recent value. In other words, monotonic-read guarantees that if a process has seen a value of x at time t, then it will never see an older value of x at a later time.

\textbf{Monotonic-Writes:} A write operation by a process on a data item x is completed before any successive write operations on x by the same process. Thus completing a write operation means that the copy on which a successive operation is preformed on reflects the effect of a previous write operation by the same process. If need be, the new write must wait for old ones to finish. This resembles FIFO consistency, in which write operations by the same process are preformed in the correct order everywhere. 

\textbf{Read-Your-Writes:} The effect of a write operation by a process on data item x will always be seen by a successive read operation on x by the same process. Write operations are always completed before a successive read operation by the same process, no matter where the read operation takes place. 

\textbf{Write-Follow-Reads:} A write operation by process on a data item x following a previous read operation on x by the same process is guaranteed to take place on the same or more recent value of x that was read. Successive write operations by a process on a data item x will be preformed on a copy of x that is up to date with the value most recently read by that process.


\end{minipage} % End the main body - first page mini page

\end{document} 
